% Part: lambda-calculus
% Chapter: lambda-definability
% Section: primitive-recursive-functions

\documentclass[../../../include/open-logic-section]{subfiles}

\begin{document}

\olfileid[id]{lam}{ldf}{prf}
\olsection{Fungsi Rekursif Primitif yang \usetoken{S}{lambda definable}}

Ingat bahwa fungsi rekursif primitif ialah fungsi yang dapat didefinisikan
dari fungsi-fungsi dasar $\Zero$, $\Succ$, dan $\Proj{n}{i}$ melalui
komposisi dan rekursi primitif.

\begin{lem}
  \ollabel{lem:basic}
  Fungsi-fungsi rekursif primitif dasar $\Zero$, $\Succ$, dan
  proyeksi~$\Proj{n}{i}$ bersifat !!{lambda definable}.
\end{lem}

\begin{proof}
Fungsi-fungsi tersebut !!{lambda defined} oleh suku-suku berikut:
\begin{align*}
  \fn{Zero} & \ident \lambd[a][\lambd[fx][x]]\\
  \fn{Succ} & \ident \lambd[a][\lambd[fx][f (a f x)]] \\
  \fn{Proj}^n_i & \ident \lambd[x_0\dots x_{n-1}][x_i]
\end{align*}
\end{proof}

\begin{lem}
  \ollabel{lem:comp} Andaikan fungsi beraritas~$k$, $f$, dan fungsi-fungsi
  beraritas~$n$, $g_0, \dots, g_{k-1}$, masing-masing
  !!{lambda defined} oleh suku $F$, $G_0$, \dots,~$G_{k-1}$, dan $h$
  didefinisikan dari fungsi-fungsi tersebut melalui komposisi. Maka $h$
  bersifat !!{lambda definable}.
\end{lem}

\begin{proof}
  Fungsi $h$ dapat !!{lambda defined} oleh suku
  \[
  H \ident \lambd[x_0 \dots x_{n-1}][F\, (G_0 x_0 \dots
    x_{n-1}) \dots (G_{k-1} x_0 \dots x_{n-1})]
  \]
  Verifikasi fakta ini diserahkan sebagai latihan.
\end{proof}

\begin{prob}
  Lengkapi bukti \olref[lam][ldf][prf]{lem:comp} dengan menunjukkan
  bahwa $H\num{n_0}\dots\num{n_{n-1}} \red \num{h(n_0, \dots,
    n_{n-1})}$.
\end{prob}

Perhatikan bahwa \olref{lem:comp} tidak mensyaratkan $f$ dan $g_0$,
\dots,~$g_{k-1}$ bersifat rekursif primitif; yang disyaratkan hanyalah bahwa
fungsi-fungsi tersebut total dan !!{lambda definable}.

\begin{lem}\ollabel{lem:prim}
  Andaikan $f$ adalah fungsi beraritas~$n$ dan~$g$ adalah fungsi
  beraritas~$n+2$, keduanya !!{lambda defined} oleh suku $F$ dan~$G$, dan
  fungsi~$h$ didefinisikan dari $f$ dan $g$ melalui rekursi primitif.
  Maka $h$ juga bersifat !!{lambda definable}.
\end{lem}

\begin{proof}
  Ingat bahwa $h$ didefinisikan oleh
  \begin{align*}
    h(x_1, \dots, x_n, 0) &= f(x_1, \dots, x_n)\\
    h(x_1, \dots, x_n, y+1) & = g(x_1, \dots, x_n, y, h(x_1, \dots, x_n, y)).
  \end{align*}
  Secara informal, definisi rekursif primitif mengiterasikan langkah yang
  ditentukan oleh fungsi $g$ sebanyak $y$ kali, dengan bertolak dari
  $f(x_1, \dots, x_n)$. Hal ini mengingatkan kita pada definisi numeral
  Church, yang juga didefinisikan sebagai iterator.

  Demi kesederhanaan, kita memberikan definisi dan bukti untuk satu argumen
  parameter tambahan~$x$. Fungsi $h$ !!{lambda defined} oleh:
  \begin{align*}
    H \ident & \lambd[x][\lambd[y][\fn{Snd} (y
        D \tuple{\num{0}, F x})]]
    \intertext{dengan}
    D \ident & \lambd[p][\tuple{\fn{Succ} (\fn{Fst}\, p), (G x
          (\fn{Fst}\, p) (\fn{Snd}\, p))}]
  \end{align*}
  Keadaan iterasi yang kita pertahankan adalah suatu pasangan: komponen
  pertamanya ialah nilai $y$ saat ini, sedangkan komponen keduanya ialah nilai
  $h$ yang bersesuaian. Pada setiap langkah iterasi, kita membuat pasangan
  nilai $y$ dan $h$ yang baru; setelah iterasi selesai, kita mengembalikan
  komponen kedua pasangan, dan itulah nilai akhir~$h$. Sekarang kita buktikan
  bahwa ini memang merupakan representasi rekursi primitif.

  Kita ingin membuktikan bahwa untuk setiap $n$ dan $m$,
  $H\,\num{n}\,\num{m} \red \num{h(n,m)}$. Untuk itu, pertama-tama kita
  tunjukkan bahwa jika $D_n \ident \Subst{D}{\num{n}}{x}$, maka
  $D_n^m \tuple{\num{0}, F\, \num{n}} \red
  \tuple{\num{m}, \num{h(n, m)}}$. Kita melanjutkan dengan induksi pada~$m$.

  Jika $m=0$, kita ingin memperoleh
  $D_n^0 \tuple{\num{0}, F\, \num{n}} \red
  \tuple{\num{0}, \num{h(n, 0)}}$. Namun,
  $D_n^0 \tuple{\num{0}, F\, \num{n}}$ tidak lain adalah
  $\tuple{\num{0}, F\, \num{n}}$. Karena $F$ !!{lambda define}s~$f$,
  suku ini tereduksi menjadi $\tuple{\num{0}, \num{f(n)}}$; dan karena
  $f(n) = h(n, 0)$, hasilnya adalah
  $\tuple{\num{0}, \num{h(n,0)}}$.

  Sekarang andaikan $D_n^m \tuple{\num{0}, F\, \num{n}} \red
  \tuple{\num{m}, \num{h(n, m)}}$. Kita ingin menunjukkan bahwa
  $D_n^{m+1} \tuple{\num{0}, F\, \num{n}} \red
  \tuple{\num{m+1}, \num{h(n, m+1)}}$.
  \begin{align*}
    D_n^{m+1} \tuple{\num{0}, F\, \num{n}}
    & \ident D_n(D_n^{m} \tuple{\num{0}, F\, \num{n}})\\
    & \red D_n\,\tuple{\num{m}, \num{h(n, m)}}
      \text{\qquad (menurut hipotesis induksi)}\\
    & \ident (\lambd[p][
      \tuple{
        \fn{Succ} (\fn{Fst}\, p), (G \, \num{n} (\fn{Fst}\, p) (\fn{Snd}\, p))
    }]) \tuple{\num{m}, \num{h(n,m)}}\\
    & \redone
    \tuple{\fn{Succ} (\fn{Fst}\, \tuple{\num{m}, \num{h(n,m)}}),\\
      & \qquad (G \, \num{n} (\fn{Fst}\, \tuple{\num{m}, \num{h(n,m)}}) (\fn{Snd}\, \tuple{\num{m}, \num{h(n,m)}}))}\\
    & \red
    \tuple{\fn{Succ}\, \num{m}, (G \, \num{n} \,\num{m} \, \num{h(n,m)})}\\
    & \red \tuple{\num{m+1}, \num{g(n, m, h(n, m))}}
  \end{align*}
  Karena $g(n, m, h(n, m)) = h(n, m+1)$, pembuktian langkah induksi selesai.

  Terakhir, tinjau
  \begin{align*}
    H\,\num n\,\num m &\ident \lambd[x][\lambd[y][\fn{Snd} (y
        (\lambd[p][\tuple{\fn{Succ} (\fn{Fst}\, p), (G\, x\,
          (\fn{Fst}\, p)\, (\fn{Snd}\, p))}]) \tuple{\num{0}, F x})]]\\
    & \qquad\,\num n\, \num m\\
    & \red \fn{Snd} (\num m \,
    \underbrace{(\lambd[p][\tuple{\fn{Succ} (\fn{Fst}\, p), (G \,\num n\,
      (\fn{Fst}\, p) (\fn{Snd}\, p))}])}_{D_n} \tuple{\num{0}, F \num n})\\
    & \ident \fn{Snd} (\num m \, D_n\, \tuple{\num{0}, F \num n})\\
    & \red \fn{Snd} \,(D_n^m  \tuple{\num{0}, F \num n})\\
    & \red \fn{Snd} \,\tuple{\num{m}, \num{h(n, m)}} \\
    & \red \num{h(n,m)}.
  \end{align*}
\end{proof}


\begin{prop}
  Setiap fungsi rekursif primitif bersifat !!{lambda definable}.
\end{prop}

\begin{proof}
  Berdasarkan \olref{lem:basic}, semua fungsi dasar bersifat
  !!{lambda definable}; dan berdasarkan \olref{lem:comp} serta
  \olref{lem:prim}, fungsi-fungsi yang !!{lambda definable} tertutup terhadap
  komposisi dan rekursi primitif.
\end{proof}

\end{document}
